--- Page 1 ---
Published:
March 18, 2011
r 2011 American Chemical Society
5594
dx.doi.org/10.1021/ja200560z | J. Am. Chem. Soc. 2011, 133, 5594–5601
ARTICLE
pubs.acs.org/JACS
Valence-to-Core X-ray Emission Spectroscopy: A Sensitive Probe of
the Nature of a Bound Ligand
Christopher J. Pollock and Serena DeBeer*
Department of Chemistry and Chemical Biology, Cornell University, Ithaca, New York 14853, United States
b
S Supporting Information
’INTRODUCTION
Metalloenzymes containing iron are of critical importance to
many biological processes, including hydrocarbon oxidation and
nitrogen fixation.1�4 Many of these reactions are of considerable
interest on both laboratory and industrial scales, fueling investi-
gation into these enzymes with the hope of mimicking their
reactivities. Harnessing and replicating these reactivities requires
detailed knowledge of the enzymes’ geometric and electronic
structures; to this end, crystallography and spectroscopy have
shed much light onto many of these systems. Metal spin and
oxidation states of resting enzymes and even many intermediates
have been determined by electron paramagnetic resonance
(EPR), X-ray absorption (XAS), X-ray magnetic circular dichro-
ism (X-MCD), and M€ossbauer, while structural information has
been obtained from resonance Raman, extended X-ray absorp-
tion fine structure (EXAFS), and magnetic circular dichroism
(MCD).5�11
Despite the insights from these tools, many questions about
these metal centers and their ligand frameworks still remain,
especially with respect to reactive intermediates. The oxygen
binding mode in intermediate Q of methane monooxygenase12
and the protonation state of the compound II intermediates in
hemeenzymes8,13providetwo examples of enduring uncertainties.
One reason for the persistence of these questions is the lack of
spectroscopic techniques that readily probe the bound ligands.
Of the techniques that are sensitive to ligand identity, EXAFS
is unable to differentiate between C, N, and O scatterers, while
EPR and vibrational analyses require isotope substitution, often a
prohibitive challenge in a protein matrix. Perhaps the most
sensitive well-established probe of ligand identity is photoemission
spectroscopy, but here the requirement for ultrahigh vacuum
conditions greatly limits the systems to which it may be applied.
Even if these aforementioned challenges could be overcome, it
is still exceedingly difficult to probe beyond the atoms immedi-
ately bound to the metal. EXAFS, for example, could potentially
differentiate between a bound water and a hydroxide ion based
on bond length differences, but may not be able to discriminate a
water from a carboxylate. While it is possible to investigate
beyond the first coordination sphere with complex ENDOR,
ESEEM, and NMR experiments, these all require the presence of
sizable couplings, limiting their application.14,15 The power to
easily distinguish between ligands that provide similar coordina-
tion environments, such as water versus oxidized substrate, is of
utmost importance in complex systems such as enzyme active
sites, although to do so requires a currently unavailable means of
probing the ligand electronic structure. Thus, the difficulties and
limitations endemic to currently available methods warrant the
development of new techniques capable of more thoroughly
investigating the ligands around a metal center.
It is here, as a spectroscopic probe of the ligand environment,
that Fe Kβ X-ray emission spectroscopy (XES) holds promise.16,17
This technique involves ionizing a 1s core electron on Fe by
incident high energy X-rays, followed by monitoring the emission
of photons during electron decay to fill the 1s hole (Figure 1).18
The features of chemical interest are the Kβ1,3 emission, an electric
dipole-allowed 3p�1s transition that is dominated by 3p�3d
exchange correlation with a lesser 3p spin�orbit coupling com-
ponent, and the Kβ00/Kβ2,5 “valence-to-core” emissions, which
stem from transitions from the filled ligand valence ns and np
orbitals, respectively, to the Fe 1s core hole.19 Fe Kβ XES has
already been shown to be sensitive both toward metal spin and
oxidation state, as well as ligand identity, hybridization, and
Received:
January 19, 2011
ABSTRACT: The sensitivity of iron Kβ X-ray emission spectros-
copy (XES) to the nature of the bound ligands (σ-donating, π-
donating, and π-accepting) has been explored. A combination of
experiment and theory has been employed, with a DFT approach
being utilized to elucidate ligand effects on the spectra and to
assign the spectral intensity mechanisms. Knowledge of the
various contributions to the spectra allows for a deeper under-
standing of spectral features and demonstrates the sensitivity of
this method to the identity of the interacting ligands. The potential
of XES for identifying intermediate species in nonheme iron
enzymes is highlighted.
Downloaded via FREIE UNIV BERLIN on November 28, 2025 at 12:32:50 (UTC).
See https://pubs.acs.org/sharingguidelines for options on how to legitimately share published articles.


--- Page 2 ---
5595
dx.doi.org/10.1021/ja200560z |J. Am. Chem. Soc. 2011, 133, 5594–5601
Journal of the American Chemical Society
ARTICLE
protonation.16 Of particular interest is the valence-to-core region
because these features represent transitions from filled orbitals that
are dominantly ligand in nature. Despite the rich information
content of this region, the exact mechanisms and selection rules
underpinning intensities have yet to be fully understood.
Fortunately, Fe Kβ XES spectra can be modeled with excellent
agreement to experiment by density functional theory (DFT),
thus providing a means to gain further insight into these spectra.
Figure 2 shows a comparison of the experimental XES data
for Fe(TACN)2
3þ, Fe(acac)3, and Fe(CN)6
3� (TACN = 1,4,7-
triazacyclononane, acac = acetylacetonate) and the correspond-
ing calculated spectra.16 Analysis of the dominant contributions
by DFT demonstrated that the valence-to-core region of XES
spectra is comprised primarily of contributions from the ligands.
As shown in Figure 2, in the case of a pure σ-donating ligand
(such as TACN), only a single feature in the XES Kβ2,5 region is
observed, while for the π-donating and π-accepting ligands, acac
and CN�, respectively, two features are observed. It thus seems
intuitive that the nature of the ligand must influence the spectra,
although a thorough investigation of ligand contributions has not
been previously undertaken. For instance, contributions from the
σ and π frameworks have not been established, nor are the
intensity mechanisms fully understood. To exploit the potential
of XES for investigating proteins, detailed knowledge of exactly
why these features arise for π-donors and acceptors but not
for σ-only donors, as well as more general knowledge about what
governs intensity, is required.
The present study employs a DFT approach to investigate in
detail the origins of the valence-to-core features and the mechan-
isms governing their intensity. As the iron valence-to-core
spectra have been previously shown to be dominated by Fe np
to Fe 1s transitions, we have undertaken a theoretical investigation
from the perspective of both the metal emission and the ligand
emission (Figure 3) to unambiguously understand the ligand
contributions.16 For first row elements, examining the KR ligand
(nitrogen or oxygen in the current investigation) emissions is
critical because they provide a measure of the ligand electronic
structure,20,21 similar to photoemission spectroscopy, that can be
correlated to the iron XES spectra, thus giving a more complete
picture of how the ligand contributes to the experimental spectra.
These studies form a foundation for establishing this technique as
a selective probe of the ligand environment and determining the
factors governing intensity. The 4p contributions of the metal to
bonding orbitals are also discussed. Finally, the extension of these
methods and the possibility for using valence-to-core XES for
understanding metal substrate interactions in metalloprotein
active sites is highlighted.
Figure 1. Fe Kβ XES spectrum of Fe2O3 with valence-to-core region
enlarged roughly 100-fold. The simplified molecular orbital (MO)
diagram on the right schematically depicts the origin of the features in
the spectrum.
Figure 2. Comparison between experimental and DFT calculated valence-to-core Fe Kβ XES spectra for Fe(TACN)2
3þ, Fe(acac)3, and Fe(CN)6
3�. A
2.5 eV broadening and 182.5 eV energy shift have been applied to the calculated spectra.16
Figure 3. Schematic of the two types of transitions employed in this
work: Fe np to Fe 1s (Fe Kβ valence-to-core XES, blue) and ligand 2p to
ligand 1s (ligand KR XES, red).


--- Page 3 ---
5596
dx.doi.org/10.1021/ja200560z |J. Am. Chem. Soc. 2011, 133, 5594–5601
Journal of the American Chemical Society
ARTICLE
’CALCULATIONS
All calculations were performed using the ORCA quantum chemistry
software package.22NitrogenKRandiron Kβ XESspectrawere calculated
on either optimized or idealized structures (specified in text) using the
protocol described previously.16 The BP86 functional23,24 was utilized
with the TZVP25 basis set on all atoms excluding Fe, which had an
expanded CP(PPP) basis set.26 Solvation was modeled using the con-
ductor-like screening model (COSMO)27 in an infinite dielectric. Geo-
metry optimizations were performed using the same level of theory.
Constructs involving octahedral arrangements of ligands around point
charges
were
obtained
by
first
optimizing
the
corresponding
ferric structure followed by subsequent replacement of the iron with a
point charge; nofurther optimizations werecarried out. Molecularorbitals
were visualized using Chimera.28 A 2.5 eV broadening was applied to all
calculated iron emission spectra, and a 1 eV broadening was applied to all
calculated ligand (N and O) emission spectra.
’RESULTS AND ANALYSIS
Mechanisms Governing Spectral Intensity. Because the
nature of the ligands bound to the metal was found to greatly
impact XES spectra,16 emphasis was placed on correlating the
ligand electronic structure to the valence-to-core region of Fe XES
spectra.For ease of analysis, a complex with only simple σ-donating
ligands was considered first (ferric hexammine, Fe(NH3)6
3þ).
DFT calculations were performed using established computational
methods (sample input file included in the Supporting In-
formation) on free NH3 without any metal present to establish
the ligand’s electronic structure. The structure was first optimized,
and then a nitrogen XES spectrum was calculated using transitions
from the valence N 2p orbitals to a N 1s hole (KR emission),
allowing the observation of the energetics of the NH3 molecular
orbitals (Figure 4). Clear transitions are observed from all expected
valencemolecularorbitals(MOs),andtherelativeenergiesobtained
are in good agreement with previous photoelectron spectroscopic
results.29�31 Furthermore, the XES calculation was used to assess
the relative contributions to the spectrum from electric dipole,
magnetic dipole, and electric quadrupole mechanisms. Analysis of
these transitions revealed they derived almost exclusively from
electric dipole-allowed character (>99.5%), as would be expected
for KR emission.
Once the features of the free NH3 emission spectrum were
established, another level of complexity was added by arranging
six ammonias into an octahedral geometry around a central point
charge; this was equivalent to replacing the iron of Fe(NH3)6
3þ
with a point charge, Q. The charge of Q was then varied
incrementally from 0 to þ3, keeping the geometry of the ammonias
constant, and nitrogen XES spectra were calculated for each of these
charges (Supporting Information). The spectra remain similar to
that from Figure 4, although, as expected, increasing the charge of Q
resulted in increasing stabilization of the 1a2 lone pair orbital, up to
roughly 3 eV, while the 1a1 and 1e N�H bonding orbitals were
slightly destabilized (<0.7 eV). Slight additional splitting is also seen
for the 1a2 feature resulting from the various phase combinations
of the six lone pairs in octahedral symmetry (a1 g vs eg vs t1u).
Using this same octahedral framework, the point charge was
replaced by an Fe3þ ion, and the N emission was again analyzed
(Figure 5). The spectrum thus obtained is remarkably similar,
both in terms of energy and intensity of transitions, to that with
only a Q3þ instead of Fe3þ, suggesting that an iron ion and a
point charge are similar perturbations on the ligand orbitals.
The identification of the spectral features in the nitrogen KR
emission spectrum of Fe(NH3)6
3þ allows for direct comparison to
the iron Kβ emission for the same compound by aligning the N 2s
emission feature in both spectra (Figure 6). Excellent energetic
agreement is seen for all features between the two spectra,
providing further proof that iron valence-to-core Kβ XES spectra
are dominated by ligand electronic structure. The relative inten-
sities, however, differ markedly between the two cases. A factor of
1000 difference in intensity is expected between the KR and Kβ2,5
emission spectra (nitrogen and iron, respectively) due to the
relative probabilities of each of these transitions,19 with additional
contributions possible due to differences in the intrinsic dipole
moment of Fe and N.
Qualitatively, it can be seen that intensity is roughly modulated
by the degree of overlap between the ligand and metal molecular
orbitals; there is less overlap between the metal and the N 2s
orbital as compared to the lone pair, and this is reflected in the
relative intensities. Closer inspection of the transitions in the iron
spectrum, however, reveals that all features with significant
intensity derive from combinations of ammonia molecular
orbitals possessing t1u symmetry (notice the unresolved eg lone
pair feature at 6931 eV). Notably, this is the same symmetry as
the iron p orbitals in octahedral geometry, implicating metal p
Figure 4. Comparison showing a nitrogen KR XES spectrum of NH3
(N 2p f N 1s transitions, left) with transitions labeled according to their
symmetry from the NH3 molecular orbital diagram (right).
Figure 5. Nitrogen emission spectra from a construct of six ammonias
around a central moiety, either a point charge of 0 or þ3 or a ferric ion.
Spectra are quite similar in energy and intensity of features, except for
stabilization of the 1a2 lone pair by charge as indicated by arrow.


--- Page 4 ---
5597
dx.doi.org/10.1021/ja200560z |J. Am. Chem. Soc. 2011, 133, 5594–5601
Journal of the American Chemical Society
ARTICLE
orbital participation in bonding that introduces electric dipole-
allowed p�s character to these transitions (vide infra).
From the N KR emission calculations, it is observed that
introducing a point charge within a ligand framework stabilizes
the ligand lone pairs without dramatically affecting intensities.
A ferric ion has an effect very similar to that of a simple þ3
point charge, stabilizing the lone pairs without impacting
transition intensities. Analysis of calculated Fe valence-to-core
Kβ emission reveals that the intensities are governed both by
overlap with the metal orbitals and also by the symmetry of
this overlap.
Given the ease with which Fe(NH3)6
3þ could be analyzed, a
similar approach was taken with the more complex σ-only donor
TACN. N KR emission spectra were calculated for Q(TACN)2
0,
Q(TACN)2
3þ, and Fe(TACN)2
3þ, and, while the MO picture
becomes more complicated, the same trends are observed that
were seen with NH3 as a ligand (Figure 7). Comparison to the Fe
emission spectrum again shows intensity modulation by overlap
between metal and ligand orbitals and by the symmetry of TACN
molecular orbitals (see the Supporting Information).
While both of these examples are insightful and help to
better understand XES spectra, they both involve only simple
σ-donors and do not explain all observed features of the spectra,
such as the split sometimes observed in the Kβ2,5 feature when π-
accepting or donating ligands are present. To delve deeper
into the mechanisms governing the spectra, a complex with π-
accepting ligands, Fe(CN)6
3�, was also studied. An MO diagram
and assigned N KR emission spectrum of CN� are shown in
Figure 8.
Not surprisingly, employing an octahedral construct of six
cyanides around a point charge similar to that used for NH3
resulted in slight destabilization of most orbitals except for the lone
pair, which was strongly stabilized (see the Supporting In-
formation). The interesting contribution to the Fe XES spectrum
is seen when it is overlaid with the N emission spectrum; the
dominantcontributions to the Kβ2,5region comefromthecyanide
σ*2s�2s and σ2p�2p orbitals (Figure 9). Moreover, there is no
significant contribution from the cyanide π system, presumably
because the filled π orbitals of CN� are localized between
the carbon and nitrogen and have poor overlap with the metal.
Figure 6. Comparison of N and Fe emission spectra from Fe(NH3)6
3þ showing intensity modulation from the perspective of iron. Representative MOs
from each transition are plotted on the right.
Figure 7. N KR XES spectra of Q(TACN)2
0, Q(TACN)2
3þ, and
Fe(TACN)2
3þ demonstrating effects similar to those seen when NH3
was used as a ligand.


--- Page 5 ---
5598
dx.doi.org/10.1021/ja200560z |J. Am. Chem. Soc. 2011, 133, 5594–5601
Journal of the American Chemical Society
ARTICLE
As can clearly be seen from the molecular orbital plots A�D in
Figure 9, intensity again derives from orbitals of t1u symmetry.
Fe(acac)3 is the only remaining compound from this series
that requires analysis. The electronic complexity of the acac
ligand obscures the clear-cut assignments used above, although
the same general trends are apparent and are given in the
Supporting Information.
Throughout this investigation, it has been seen that only
transitions from molecular orbitals of t1u symmetry show appre-
ciable intensity in calculated Fe Kβ XES spectra. In octahedral
symmetry, the metal p orbitals also transform as t1u, allowing
them to interact with any ligand-based orbitals with the same
symmetry. It is this metal p contribution that imparts electric
dipole-allowed character to MOs of t1u symmetry, providing a
mechanism for intensity not present for MOs of other symme-
tries. Furthermore, the interaction with metal p orbitals is
mediated by distance as well, explaining the higher intensity
typically seen for low spin complexes over high spin ones.
Metal p Contributions to Bonding. To quantify the extent of
metal p contributions to bonding, a correlation was sought between
the amount of metal p character in a molecular orbital and the
intensity of the corresponding XES transition signal. For all four
compounds studied here, the oscillator strength of all transitions, a
sum of electric dipole, magnetic dipole, and quadrupole contribu-
tions, was plotted against the amount of metal p character from each
transition (Figure 10). The strong correlation obtained suggests
that metal p character mixed into these orbitals provides the
mechanism for intensity, as has been previously suggested.16,32,33
No correlation was found between oscillator strength and metal s,
ligand s, or ligand p character, and no stronger correlation was seen
when only electric dipole moments were considered (see the
Supporting Information). Additionally, the quadrupole contribu-
tion to the spectra was negligible, comprising less than 3% of
observed intensity, similar to what has been observed previously.16
With p-character established as contributing to intensity, the
next question became the origin of this p contribution: is it from
metal 3p or 4p orbitals? Previous attempts to determine the
nature of the metal np mixing employed the atomic natural
orbital basis set ano-pVDZ and showed that both 3p and 4p
contributions to the valence-to-core spectra were present.16 This
study indicated that, on one hand, Fe 3p character is borrowed
from the Kβ main line and, on the other hand, 4p character may
contribute to bonding. Here, we take a slightly different approach
and instead ask if there is there any evidence for 4p contributions
to the filled bonding orbitals.
From the output of the DFT calculations considered here, the
total Fe p contribution to filled orbitals was calculated for each of
Figure 9. Correlation between Fe and N emission spectra of Fe(CN)6
3� showing the origin of features in Fe spectrum. The split Kβ2,5 peak results from
σ interactions between the metal and the σ*2s�2s and σ2p�2p orbitals on the cyanide ligands (as seen from representative MOs B and C, respectively,
plotted on the right). No significant contribution from the π system is observed.
Figure 8. Molecular orbital diagram and N XES spectrum of CN�.


--- Page 6 ---
5599
dx.doi.org/10.1021/ja200560z |J. Am. Chem. Soc. 2011, 133, 5594–5601
Journal of the American Chemical Society
ARTICLE
the four complexes studied (Table 1). In all cases, the total Fe p
character exceeded the 1200% expected for completely filled 2p
and 3p orbitals (300% for each spin up and spin down electrons
in both 2p and 3p), necessitating the invocation of 4p contribu-
tions to the filled MOs. The 4p values were obtained by simply
subtracting 1200% from the total p amount to account for filled
2p and 3p orbitals. Furthermore, examination of the Loewdin
orbital populations revealed negligible ligand contributions to
the metal 3p orbitals (<1%), confirming that the p contribution
here derives from the 4p and not the 3p orbitals. This supports
the notion that there is very little overlap between the metal 3p
and the valence ligand orbitals.
For all compounds, the amountofp character is roughly 0.8�1.5
p electrons, although what this is worth in terms of bond strength is
still uncertain. To ensure that the presence of p character beyond
1200% was not simply an artifact of the calculation, the electronic
structure of a free ferric ion was performed, confirmingfilled 2p and
3p and empty 4p, as indicated in Table 1.
’DISCUSSION
The data presented herein demonstrate that valence-to-core
Fe Kβ XES spectra are dominated by the electronic structure of
the ligand, to the extent that the spectra can be interpreted as
direct probes of the ligand molecular orbitals. These ligand
orbitals are stabilized to a significant degree by the charge on
the metal, thus determining the energy of the transitions seen in
the XES spectrum, while the intensity is governed by the extent of
Fe 4p mixing. The latter point is of particular interest as there
now exists a means to experimentally quantify the extent of 4p
contributions to bonding, which could lead to new insights into
bonding in transition metal complexes, a topic that has been the
subject of some controversy.34�36 The contributions of both the
ligand MOs and the Fe 4p interaction combine to make valence-
to-core Fe XES a uniquely sensitive technique for investigating
Fe coordination in complex environments.
Of particular promise for XES is its application to active site
intermediates in iron-containing enzymes. In many cases,
mechanistic proposals present competing pathways that are
difficult to distinguish using traditional chemical methods. One
such example is the extradiol dioxygenase family of enzymes
that catalyze the oxidative ring cleavage of catechol derivatives.
For many years, the intermediates along the reaction cycle were
debated (Figure 11); specifically, it was unclear whether a side-on
or end-on superoxide is formed and also whether a lactone or
dioxetane intermediate is involved.1 Recently, these questions
were largely answered by a fortuitous crystal structure of one
enzyme from this family,37 although occurrences such as this are
surely the exception rather than the rule. Had a crystal structure
not been available, as is the case for most reactive intermediates
and many other members of this enzyme family, the questions
surrounding this enzyme are such that XES would likely be able
to differentiate between the competing possibilities. This method
thus holds promise for characterization of other members of this
family of enzymes, and broadly for iron enzyme active site/
substrate interactions in general.4,5,38,39 Of particular interest is
the application of XES to reactive intermediates that are not
amenable to current structural characterization techniques.
For XES analysis, the two possible orientations of the super-
oxide should be distinguishable due to differing levels of stabi-
lization and Fe 4p mixing for the various orbitals interacting with
the metal. The side-on conformation would be expected to
interact strongly through the 2p�2p π bonding orbital with a
weaker interaction between the 2p�2p π*; the end-on species, in
contrast, would bond strongest through 2p�2p π* orbitals with
the 2p�2p π forming a secondary interaction. The differential
stabilization of, and Fe 4p contributions to, these orbitals would
Table 1. Metal 4p Character in the Compounds Studied
compound
total % p
% 4p
Fe3þ
1199.6%
0%
Fe(NH3)6
3þ
1280.6%
80.6%
Fe(CN)6
3�
1345.1%
145.1%
Fe(TACN)2
3þ
1295.5%
95.5%
Fe(acac)3
3�
1284.4%
84.4%
Figure 11. A condensed catalytic cycle of the extradiol dioxygenase family of enzymes showing competing proposed intermediates: a side-on superoxide
(1a) versus an end-on superoxide (1b) and a lactone (2a) versus a dioxetane (2b). Figure adapted from ref1.
Figure 10. A strong correlation (R = 0.93) is seen between the amount
of p character of a given transition and the oscillator strength. The small
cluster of points between 2% and 7% p character that appears to fall at
lower than expected intensity consists of mostly Kβ00 transitions; the
reason for this low intensity is not yet understood.


--- Page 7 ---
5600
dx.doi.org/10.1021/ja200560z |J. Am. Chem. Soc. 2011, 133, 5594–5601
Journal of the American Chemical Society
ARTICLE
lead to valence-to-core XES features with characteristic energies
and intensities that could be used to distinguish between the two
orientations.
Similarly, the lactone and dioxetane would also possess
different electronic structures and metal�ligand interactions.
In the case of the dioxetane, the Fe is bound by two alkoxide
oxygens, while the lactone has one alkoxide, a ketone carbonyl,
and a hydroxide anion as ligands. The principal difference
between these two frameworks is the presence of the carbonyl
in the dioxetane; this functional group containing an sp2 hybri-
dized oxygen would be expected to give rise to significantly
different features from the sp3 hybridized oxygens of the dioxe-
tane. For example, the C�O 2s�2s σ bonding orbital that can
interact with the Fe would be expected to be of lower energy for
the sp2 hybridized ketone as compared to the sp3 hybridized
alkoxides, providing one means to differentiate the lactone from
the dioxetane. Perhaps more significantly, the presence of the
CdO would also result in significantly greater stabilization of the
C�O 2pz�2pz σ bonding orbital, pulling this feature to lower
energy as compared to the corresponding alkoxide and possibly
even splitting it off from the main peak (similar to what was seen
for CN�, vide supra). Because of these effects, it is expected that
XES would be able to distinguish these two intermediates as well
(Figure 12).
’CONCLUSIONS
A systematic computational study of the mechanisms govern-
ing valence-to-core Fe Kβ XES spectra has been presented. The
features seen in these spectra derive from ligand orbitals and are
dominated by two factors: stabilization due to the charge on the
metal and overlap with the metal 4p orbitals. Positive charge on
the metal significantly stabilizes the donor ligand orbitals to lower
energy (roughly 2�3 eV for ferric ions), and thus the transitions
seen in valence-to-core Fe Kβ XES are at lower energy than
would be expected for the free ligands. The stabilization observed
is very similar to that predicted for an unshielded þ3 point
charge. Intensity, meanwhile, is dictated by mixing with metal 4p
orbitals, this confers dipole-allowed character to the transitions,
which in turn is controlled by symmetry and orbital overlap.
Only ligand orbitals possessing the appropriate symmetries are
allowed to interact with the metal 4p orbitals (t1u in Oh
symmetry), so only these orbitals have the potential for signifi-
cant intensity; orbitals that cannot mix with the metal 4p give rise
to features with negligible intensity. The amount of 4p character
mixed into the ligand orbitals is governed by overlap, so short
bond lengths and ligand orbitals directed at the metal give rise to
large intensities.
Knowing the factors behind valence-to-core XES spectra
allows for better application to complicated chemical and
biological systems. Because the spectra are dependent on the
electronic structure, and thus complexity, of the ligands bound to
the metal, this technique holds promise to elucidate catalytic
intermediates, both small molecule and enzyme, during reaction
cycles that would be difficult or impossible to obtain with current
methods. This knowledge is also sufficiently general that it can be
applied to other metals, greatly broadening the scope of valence-
to-core XES spectra.
’ASSOCIATED CONTENT
b
S
Supporting Information.
Sample ORCA input file; cal-
culated N emission spectra for Q(NH3)6
nþ (n = 0�3); compar-
ison between N and Fe emission for Fe(TACN)2
3þ; calculated
N emission for Q(CN)6
n� (n = 3, 6); calculated O emission for
acac constructs; correlations between oscillator strength and
metal s character, ligand s character, and ligand p character;
and correlation between dipole-only oscillator strength and
metal p character. This material is available free of charge via
the Internet at http://pubs.acs.org.
’AUTHOR INFORMATION
Corresponding Author
serena.debeer@cornell.edu
’ACKNOWLEDGMENT
S.D. thanks Cornell University and the American Chemical
Society Petroleum Research Fund (50270-DNI3) for funding.
’REFERENCES
(1) Solomon, E. I.; Brunold, T. C.; Davis, M. I.; Kemsley, J. N.; Lee,
S. K.; Lehnert, N.; Neese, F.; Skulan, A. J.; Yang, Y. S.; Zhou, J. Chem. Rev.
2000, 100, 235–349.
(2) Baik, M. H.; Newcomb, M.; Friesner, R. A.; Lippard, S. J. Chem.
Rev. 2003, 103, 2385–2419.
(3) Costas, M.; Mehn, M. P.; Jensen, M. P.; Que, L. Chem. Rev. 2004,
104, 939–986.
(4) Hoffman, B. M.; Dean, D. R.; Seefeldt, L. C. Acc. Chem. Res. 2009,
42, 609–619.
(5) Bollinger, J. M.; Diao, Y.; Matthews, M. L.; Xing, G.; Krebs, C.
Dalton Trans. 2009, 905–914.
(6) Sarangi, R.; Gorelsky, S. I.; Basumallick, L.; Hwang, H. J.; Pratt,
R. C.; Stack, T. D. P.; Lu, Y.; Hodgson, K. O.; Hedman, B.; Solomon, E. I.
J. Am. Chem. Soc. 2008, 130, 3866–3877.
(7) Chow, M. S.; Eser, B. E.; Wilson, S. A.; Hodgson, K. O.; Hedman,
B.; Fitzpatrick, P. F.; Solomon, E. I. J. Am. Chem. Soc. 2009, 131,
7685–7698.
(8) Green, M. T.; Dawson, J. H.; Gray, H. B. Science 2004, 304,
1653–1656.
Figure 12. Calculated spectra for hypothetical FeIII(NH3)4X2 com-
plexes, where X = �OMe (red) or CH2O (blue), to demonstrate the
effect of oxygen hybridization on the XES spectra. The hypothetical
Fe(NH3)4(OMe)2
þ complex is shown with the spectra.




* xyz 3 6
Fe 0.00000000 0.00000000 0.00000000
N -2.00000000 0.00000000 0.00000000
H -2.42050000 0.00220000 -0.93680000
H -2.42190000 -0.81130000 0.46830000
H -2.42100000 0.81130000 0.46950000
N 2.00000000 0.00000000 0.00000000
H 2.42050000 -0.00220000 -0.93680000
H 2.42190000 0.81130000 0.46830000
H 2.42100000 -0.81130000 0.46950000
N 0.00000000 2.00000000 0.00000000
H -0.00220000 2.42050000 -0.93680000
H 0.81130000 2.42190000 0.46830000
H -0.81130000 2.42100000 0.46950000
N 0.00000000 -2.00000000 0.00000000
H -0.00220000 -2.42050000 -0.93680000
H 0.81130000 -2.42190000 0.46830000
H -0.81130000 -2.42100000 0.46950000
N 0.00000000 0.00000000 -2.00000000
H -0.00220000 -0.93680000 -2.42050000
H 0.81130000 0.46830000 -2.42190000
H -0.81130000 0.46950000 -2.42100000
N 0.00000000 0.00000000 2.00000000
H -0.00220000 -0.93680000 2.42050000
H 0.81130000 0.46830000 2.42190000
H -0.81130000 0.46950000 2.42100000

